% ЕМТ 2024, VIII клас — точен LaTeX препис от официалните файлове в Klasirane.com.
% Условия: https://klasirane.com/api/competitions/EMT/All/2024/8%20%D0%BA%D0%BB/probs
% SHA-256 (условия): 9fa8f20f0099dce96a087dbf56c773d0f197bb872f1ada28be67e991f60eed68
% Решения: https://klasirane.com/api/competitions/EMT/All/2024/8%20%D0%BA%D0%BB/sol
% SHA-256 (решения): e7bd118ae762b336a124f09020f060a70f431bf37e1176d0a2f1e06b1d0c949d
% Точковите оценки, правописът, формулировките и видимите печатни особености са запазени от източника.
% И двата PDF файла съдържат сканирани страници; видимите страници са проверени пряко.

Задача 8.1. Дадени са изразите $A=x^3+2x^2y+2xy+4y^2$ и $B=x^3+3xy^2+3x^2y+2y^3$.

а) Да се разложи на два неконстантни множителя с рационални коефициенти всеки от изразите $A$, $B$ и $A+4.B$.

б) Ако $x+2y=5$ и $y+2x^2=7$, то намерете най-големия прост делител на цялото число $A+4.B$.

Решение. а) $A=x^2(x+2y)+2y(x+2y)=(x^2+2y)(x+2y)$.

$B=x^3+2x^2y+x^2y+2xy^2+xy^2+2y^3=x^2(x+2y)+xy(x+2y)+y^2(x+2y)=(x^2+xy+y^2)(x+2y)$.

$A+4.B=(x+2y)(x^2+2y+4x^2+4xy+4y^2)=(x+2y)(5x^2+4xy+4y^2+2y)$.

б) $A+4.B=(x+2y)(x^2+4xy+4y^2+2y+4x^2)=(x+2y)((x+2y)^2+2(y+2x^2))=5(5^2+2.7)=195$, чийто най-голям прост делител е 13.

Коментар. С повече усилия същият извод се получава с намиране на двойките $(x;y)$, изпълняващи даденото условие, а именно
$$
\left(\frac18(1\pm\sqrt{145});\frac1{16}(39\mp\sqrt{145})\right),
$$
и заместването им в израза, като в този случай аргументацията трябва да е валидна и за двете двойки.

Оценяване. (6 точки) 3 т. за а) (по 1 т. за вярно разлагане на всеки от изразите), 3 т. за б), от които 2 т. за изразяване на $A+4.B$ чрез $x+2y$ и $y+2x^2$ и 1 т. за правилно извършване на пресмятанията. При подход б) с намиране на $x$ и $y$: 1 т. за намиране на двете възможности за $x$ и $y$ и по 1 т. за правилно извършване на пресмятанията във всеки от двата случая.

Задача 8.2. На лист хартия е начертан правоъгълен триъгълник $ABC$ с $\angle ACB=90^\circ$ и $\angle ABC=30^\circ$. Известно е, че могат да се начертаят два кръга с радиус 1 cm върху листа така, че всяка точка от вътрешността или обиколката на триъгълника $ABC$ да лежи във вътрешността или обиколката на поне един от кръговете.

а) Покажете един възможен начин за това при $BC=3$ cm.

б) Докажете, че $BC\leq3$ cm.

Решение. а) Нека средата на $AB$ е $M$ и симетралата на $AB$ пресича $BC$ в $T$. Тогава $\angle TAB=\angle TBA=30^\circ$ и $\angle CAT=\angle BAC-\angle BAT=30^\circ$, значи $CT=\dfrac{AT}{2}=\dfrac{BT}{2}$, съответно $AT=BT=\dfrac{2BC}{3}=2$ cm и $CT=1$ cm. Сега ако $K$ и $L$ са средите на $AT$ и $BT$, то от съображения за медиана към хипотенуза следва $AK=KC=KM=KT=1$ cm и $BL=LM=LT=1$ cm. Следователно кръговете с центрове $K$ и $L$ и радиус 1 cm покриват $ABC$.

б) Нека допуснем противното. Както в а) въвеждаме $T$ и получаваме $AT=BT=\dfrac{2BC}{3}>2$ cm, а също $AB>BC>2$ cm. Така $A$, $B$, $T$ са три точки, никои две от които не лежат в кръг с радиус 1 cm (диаметър 2 cm) и значи няма как $ABC$ да е покрит от два кръга.

Коментар. Може да се докаже, че измежду всички правоъгълни триъгълници, които могат да се покрият от два кръга с радиус 1 cm, с максимално лице е именно този с катет 3 cm и прилежащ към него ъгъл $30^\circ$.

Оценяване. (6 точки) 3 т. за а), от които 1 т. за описание на центровете на двете окръжности и 2 т. за проверка, че те вършат работа; 3 т. за б), от които 1 т. за описание на три точки, всеки две от които са на разстояние над 2 cm и 2 т. за доказателство, че това е така.

Задача 8.3. Да се намерят всички естествени числа $n$, такива че
$$
a+b+c\quad\text{дели}\quad a^{2n}+b^{2n}+c^{2n}-n(a^2b^2+b^2c^2+c^2a^2)
$$
за всеки три различни естествени числа $a$, $b$ и $c$.

Отговор. $n=2$

Решение. Да изберем $b=2$, $c=1$ и нека $a\geq3$ е произволно – тогава $a+3$ дели $a^{2n}+2^{2n}+1-n(5a^2+4)$. Да положим $d=a+3$ – значи $d$ дели $(d-3)^{2n}+2^{2n}+1-n(5(d-3)^2+4)$ и значи дели и $(-3)^{2n}+2^{2n}+1-49n=9^n+4^n-49n+1$. Така числото $9^n+4^n-49n+1$ има безбройно много естествени делители и трябва непременно да е равно на 0. Директно се проверява, че това не е така за $n=1$, а при $n\geq3$ индуктивно имаме $9^n>49n$, понеже $9^3=729>441=49.9$ и (при индукционното предположение) $9^{n+1}=9.9^n>9.49n=441n>49n+49=49(n+1)$.

Остава $n=2$, което е решение за всякакви $a$, $b$, $c$, понеже
$$
a^4+b^4+c^4-2(a^2b^2+b^2c^2+a^2c^2)=(a+b+c)(a-b-c)(b-c-a)(c-a-b).
$$
(Алтернативно, $c\equiv-a-b\pmod{a+b+c}$ и значи
$$
a^4+b^4+c^4-2(a^2b^2+b^2c^2+a^2c^2)\equiv a^4+b^4+(-a-b)^4-2(a^2b^2+b^2(-a-b)^2+(-a-b)^2a^2)\pmod{a+b+c},
$$
а след разкриване на скобите се вижда, че последното е тъждествено равно на 0.)

Оценяване. (7 точки) 1 т. за разглеждане на числени стойности, по-малки или равни на 3, на две от променливите, 2 т. за достигане до зависимост с делимо, зависещо само от $n$, 1 т. за обосновка, че това делимо трябва да е 0, 1 т. за довършване на $n\neq2$ (напр. чрез неравенства с индукция), 2 т. за доказване, че $n=2$ работи.

Задача 8.4. Дадено е естествено число $n$. Равностранен триъгълник със страна $n$ е разделен на равностранни триъгълничета със страна 1; техните върхове ще наричаме възли. Равностранен триъгълник с върхове три от възлите (и страни не непременно успоредни на страните на началния) ще наричаме важен. Означаваме с $p_k$ броя ненаредени двойки различни възли, които са върхове на точно $k$ важни триъгълника. Запишете като многочлен на променливата $n$ в нормален вид изразите: а) $p_0+p_1+p_2$; б) $p_1+2p_2$.

Отговор. а), б) $\dfrac18(n^4+6n^3+11n^2+6n)$

Решение. а) Броят на възлите е $1+2+\cdots+(n+1)=\dfrac12(n+1)(n+2)$. Всяка двойка възли участва в 0, 1 или 2 специални триъгълника, така че $p_0+p_1+p_2$ е всъщност броят на всички ненаредени двойки възли:
$$
p_0+p_1+p_2=\frac12.\frac12(n+1)(n+2).\left(\frac12(n+1)(n+2)-1\right)=\frac18(n^4+6n^3+11n^2+6n).
$$

б) Имаме $p_1+2p_2=3S$, където $S$ е броят важни триъгълници, понеже всеки триъгълник е броен по веднъж откъм трите си страни, така че ще търсим $S$. Ще казваме, че важният триъгълник $\tau$ е базов, ако страните му са успоредни на тези на триъгълника със страна $n$ и е със същата ориентация. Всеки важен триъгълник $\tau$ може да се потопи в единствен базов триъгълник $f(\tau)$, чиито страни минават през върховете на $\tau$. Ако страната на базов триъгълник е равна на $k\in\{1,\ldots,n\}$ (броят на тези базови триъгълници е $1+2+\cdots+(n+1-k)=\binom{n+2-k}{2}$), то той е равен на $f(\tau)$ за $k$ различни $\tau$ (по един за всеки от възлите в основата на $f(\tau)$ без най-десния). Следователно
$$
S=\sum_{k=1}^n k\binom{n+2-k}{2}=\binom{n+3}{4}.
$$
Можем да се уверим в последното равенство така: $\binom{n+3}{4}$ е броят на всички думи с 4 „А“ и $n-1$ „Б“. Ако между крайните „А“ има $n+2-k$ букви ($k\in\{1;\ldots;n\}$), то за местата на крайните „А“ има $k$ варианта, при всеки от които за местата на средните „А“ има $\binom{n+2-k}{2}$ варианта.

(Алтернативно, използвайте без доказателство известните формули $\sum_{k=1}^n k=\dfrac{n(n+1)}2$, $\sum_{k=1}^n k^2=\dfrac{n(n+1)(2n+1)}6$ и $\sum_{k=1}^n k^3=\dfrac{n^2(n+1)^2}4$. При използването на други формули за суми следва те да бъдат доказвани.) Окончателно
$$
p_1+2p_2=3S=3\binom{n+3}{4}=\frac18(n^4+6n^3+11n^2+6n).
$$

Коментар. От а) и б) следва $p_2-p_0=(p_1+2p_2)-(p_0+p_1+p_2)=0$, т.е. $p_0=p_2$. Не ни е известно директно комбинаторно доказателство (напр. със съответствие между двойките от единия тип и двойките от другия) на този факт.

Оценяване. (7 точки) 2 т. за а), от които 1 т. за пресмятане на броя двойки възли и 1 т. за обосновка защо това е търсеният; 5 т. за б), от които 1 т. за свеждане до пресмятане на $S$, 2 т. за съответствието между $\binom{n+2-k}{2}$ базови триъгълника и $k$ важни за всяко $k$ и 2 т. за пресмятане на $\sum_{k=1}^n k\binom{n+2-k}{2}$.
